\SourceLeaf{19}
in elementa rectangula divisum: huiusmodi elementum, inter puncta quorum coordinatae sunt
\[
\begin{array}{ll}
p,&q\\
p+\mathrm{d}p,&q\\
p,&q+\mathrm{d}q\\
p+\mathrm{d}p,&q+\mathrm{d}q
\end{array}
\]
contentum, erit $=\mathrm{d}p.\mathrm{d}q$, respondebitque elemento parallelogrammatico in superficie corporis contento inter quatuor puncta, quorum coordinatae resp. erunt
\[
\begin{array}{rlll}
\mathrm{I.}&x,&y,&z\\
\mathrm{II.}&x+\lambda\mathrm{d}p,&y+\mu\mathrm{d}p,&z+\nu\mathrm{d}p\\
\mathrm{III.}&x+\lambda'\mathrm{d}q,&y+\mu'\mathrm{d}q,&z+\nu'\mathrm{d}q\\
\mathrm{IV.}&x+\lambda\mathrm{d}p+\lambda'\mathrm{d}q,&y+\mu\mathrm{d}p+\mu'\mathrm{d}q,&z+\nu\mathrm{d}p+\nu'\mathrm{d}q
\end{array}
\]
si supponimus, esse
\[
\begin{aligned}
\mathrm{d}x&=\lambda\mathrm{d}p+\lambda'\mathrm{d}q\\
\mathrm{d}y&=\mu\mathrm{d}p+\mu'\mathrm{d}q\\
\mathrm{d}z&=\nu\mathrm{d}p+\nu'\mathrm{d}q
\end{aligned}
\]
Proiectiones huius areae, quam statuimus $=\mathrm{d}s$, in tria plana axibus coordinatarum $x,y,z$ normalia, facile inveniuntur resp. $=$
\[
\begin{aligned}
&\pm(\mu\nu'-\nu\mu')\mathrm{d}p.\mathrm{d}q\\
&\pm(\nu\lambda'-\lambda\nu')\mathrm{d}p.\mathrm{d}q\\
&\pm(\lambda\mu'-\mu\lambda')\mathrm{d}p.\mathrm{d}q
\end{aligned}
\]
unde per theorema satis notum ipsa elementi area erit
\[
=\mathrm{d}p.\mathrm{d}q.\BareRoot{((\mu\nu'-\nu\mu')^2+(\nu\lambda'-\lambda\nu')^2+(\lambda\mu'-\mu\lambda')^2)}
\]
\par\endgroup
\par
Hinc patet, singula integralia in sex nostris theorematibus prolata, ad formam talem reduci $\int S\mathrm{d}p.\mathrm{d}q$, ubi $S$ vel explicite vel implicite sit functio duarum indeterminatarum $p,q$, integrationemque vel per totum planum infinitum extendendam esse, vel per eam plani partem, quae superficiem integram corporis nostri quasi repraesentat. Integratio ipsa autem modo his modo illis artificiis absolvetur, de quibus regulae generales dari nequeunt.
